\section{Introduction}
Suppose an agent is asked: \emph{What was the latest value known by the end of April?} A May release is highly relevant to the topic, easy to retrieve, and wrong for the question. On the next period the numbers change, but the required operation does not: filter evidence by the information cutoff. Similar recurring operations arise when an agent must distinguish a publication date from the reporting period it describes, or ground a numerical change in documentary context that was actually available at the time. These are not one-off facts to memorize. They are procedures that should remain useful as institutional releases evolve.

Reusable agent skills are one appealing way to persist such procedures, but the scientific object is the operation rather than the container that stores it. Demonstrating that the operation helps is subtler than comparing a targeted helper with a generic one. If generic assistance distracts the model, then the targeted helper can win while doing no better than the original agent. In that case a two-arm experiment reports a positive ``skill effect'' even though nothing has been repaired. This is not hypothetical in our data: for one frozen cutoff cell, the targeted procedure scores 100\%, the matched generic helper 72\%, and the no-skill agent 100\%. The targeted-minus-generic contrast is +28 points, but the correct repair verdict is null.

This example separates three questions that are often conflated. \emph{Repair value} asks whether T improves on the original agent N. \emph{Operation specificity} asks whether T beats structured assistance that lacks the target operation: the executed G is target-blind, while a behaviorally neutral G$_0$ would be a stronger future control. \emph{Container value} asks whether the same operation works better as callable skill state than as an operation-matched retrieval-side treatment R. The current evidence addresses the first two only under the executed N/G/T design; G$_0$ and R require new treatments and remain outside the claim.

Our question is therefore: \emph{when does a reusable temporal procedure become a binding repair for a recurring failure?} We call such a case an \emph{intervention-defined causal bottleneck}. For the same frozen endpoint we evaluate three arms: a mechanism-specific targeted procedure (T), a complexity-matched target-blind helper (G), and the original no-skill agent (N). A model--task cell qualifies only when T improves over both G and N under the family outcome frozen before model results are opened. The term \emph{bottleneck} refers only to an operation whose insertion is binding under this executed intervention; it does not claim causal exclusivity, behavioral neutrality of G, or necessity of the callable-skill container.

This definition yields three falsifiable predictions. First, T$>$G without T$>$N must be rejected as generic-helper degradation rather than counted as repair. Second, a real procedure bottleneck should be visible only where the original model has headroom; cells already at no-skill ceiling should show no identifiable gain. Third, if the reusable operation rather than endpoint-specific refitting is responsible, byte-identical skill code should retain its direction on compatible held-out releases, while transfer to an incompatible domain may legitimately be null.

The question is deliberately narrower than recent work on skill learning. Paired trace auditing, no-skill attribution, randomized skill masking, continual skill evaluation, and skill compilation already establish that reusable skills can change agent behavior and that their effects are heterogeneous~\citep{zhou2026counterfactual,dong2026agentskills,wang2026notallskills,guan2026continualskillbench,wei2026evoharness,zheng2026skillprox,xu2026hyperskill}. Likewise, RAG systems and temporal retrieval methods can implement time-aware evidence selection upstream~\citep{asai2024selfrag,yan2024crag,han2025temporalgraphs,wang2026iarag}. We therefore do not claim that reusable procedures, temporal filtering, or paired skill causality are new. Our object is the stricter \emph{diagnosis}: whether a predeclared recurring failure is actually repaired by a particular reusable operation rather than merely looking better against a weakened reference.

We study this object using source-native replay built from first-party U.S. Bureau of Economic Analysis (BEA), NOAA Climate Prediction Center, and Energy Information Administration (EIA) release histories~\citep{bea2026archive,noaa2026enso,eia2026wpsrarchive}. The study contains 35 frozen endpoints and 1,326 retained runtime-valid rows. The initial BEA/NOAA replay preserves three mechanism families motivated by TimeSage-EV---cutoff validity, release alignment, and exogenous grounding---while remaining a separate, content-addressed experimental object~\citep{yao2026timesage}. A post-review EIA expansion adds 12 cutoff endpoints using the same targeted and generic implementations byte-for-byte. Repeats test execution robustness; endpoints, not repeats, are the independent scientific units.

The results follow the three predictions rather than a universal ``skills help'' story. The no-skill anchor first overturns the apparent +28-point DeepSeek cutoff repair described above. In non-ceiling cells, however, targeted procedures beat both controls: DeepSeek documentary grounding reaches 60\% versus 20\%/20\%; held-out EIA cutoff filtering reaches 100\% versus 70\%/47.5\%, with an endpoint sign test of $p=0.0156$ for T--N; and held-out Kimi alignment reaches 12/12 versus 4/12/7/12. Transfer is conditional rather than universal: compatible future-release tests retain the expected direction, cross-domain grounding is null, and the mini model is already at no-skill ceiling on the same held-out EIA cutoff endpoints.

\paragraph{Contributions.}
Our contribution is an identification contract for reusable-procedure repair, not a new skill optimizer. We (i) define a three-arm procedure-bottleneck estimand that requires improvement over both generic assistance and the original agent; (ii) show an observed two-arm false positive in which the no-skill anchor changes the scientific verdict; and (iii) test the same frozen procedures across future releases while preserving ceiling and incompatible-domain nulls. We additionally make the identification boundary explicit: the executed contrasts estimate operation value under the frozen harness, whereas a behaviorally neutral helper and an operation-matched retrieval treatment are separate, unexecuted estimands. The resulting claim is intentionally conditional on procedure$\times$model$\times$task regime.
